\chapter{What Is LiDAR?}

\begin{kid}
Close your eyes and clap your hands in a big gym. You hear an echo. If the wall is far away, the echo takes longer to come back. If it is close, the echo comes back fast. A bat flies in the dark this way!

A LiDAR does the same thing --- but instead of a clap it uses a tiny flash of laser light, and instead of ears it has a super-fast light catcher. It does this hundreds of thousands of times every second, in every direction, so it can ``see'' the shape of everything around it: roads, trees, buildings, people. Each echo becomes one dot. Millions of dots together make a 3-D picture called a \textbf{point cloud}.
\end{kid}

\section{The name}
\textbf{LiDAR} stands for \textbf{Li}ght \textbf{D}etection \textbf{A}nd \textbf{R}anging. ``Ranging'' is the surveyor's word for measuring distance. A LiDAR is, at heart, an electronic distance meter (EDM) --- the same physics as the EDM inside a total station --- that fires extremely fast and sweeps its beam in many directions.

\fig[0.95]{pointcloud}{A synthetic mobile-LiDAR point cloud of a street, colored by height: road, curbs, poles, trees and building faces. Every dot is one laser echo.}

\section{The four families of LiDAR}
\fig[1.0]{family}{The same measuring idea mounted four ways.}

\begin{center}\small
\renewcommand{\arraystretch}{1.2}
\begin{tabularx}{\linewidth}{lXXl}
\toprule
\textbf{Platform} & \textbf{How it moves} & \textbf{Where the position comes from} & \textbf{Typical use}\\ \midrule
Airborne & aircraft / drone & GNSS + IMU (very open sky) & terrain, floodplains \\
Terrestrial (TLS) & stands still on a tripod & surveyed setups, targets & buildings, plants, forensics\\
Mobile (MLS) & car, train, boat & GNSS + IMU + LiDAR SLAM & corridors, roads, rail\\
Backpack / handheld & a walking person & IMU + LiDAR SLAM (+ control) & indoors, tunnels, under canopy\\
\bottomrule
\end{tabularx}
\end{center}

This course concentrates on the two moving families at the bottom of the table: \textbf{mobile} and \textbf{handheld} systems. They are the hardest, because the instrument never holds still --- and the most productive, because you can capture a kilometer of street in the time it takes to set up a total station a few times.

\section{The one big problem: the instrument is moving}
\begin{bigidea}
A laser can only tell you \emph{``that wall is 25.000 m away in that direction from me.''} If you do not know exactly where ``me'' is --- and which way ``me'' was pointing --- at the \emph{exact instant} of each pulse, the measurement is useless. Most of this course is about knowing ``me'': the sensor's \textbf{pose}.
\end{bigidea}

\fig[0.95]{birdsnest}{Left: the same scans placed with a bad trajectory form the famous ``bird's nest.'' Right: with a correct high-rate trajectory every scan stacks into a crisp map.}

\begin{eng}
For a moving scanner each laser return must be placed in the project frame using the full six-degree-of-freedom pose (three positions $X,Y,Z$ and three angles roll, pitch, yaw) of the sensor at the time that particular pulse was fired. A spinning scanner at 10\,Hz fires up to a few million pulses per second, so the pose must be known at very high rate. That is achieved by combining:
\begin{enumerate}
\item an \textbf{inertial measurement unit (IMU)} for high-rate, short-term motion (Module 7);
\item \textbf{precise time synchronization} so every measurement shares one clock (Module 8);
\item measured \textbf{lever arm and boresight} between the sensors (Module 9);
\item a \textbf{trajectory estimator} --- GNSS/INS fusion and/or LiDAR SLAM --- that keeps drift bounded (Modules 10--12);
\item \textbf{ground control and independent checkpoints} that tie the result to a datum and prove its accuracy (Module 13).
\end{enumerate}
\end{eng}

\section{Why surveyors already understand most of this}
\begin{center}\small
\renewcommand{\arraystretch}{1.2}
\begin{tabularx}{\linewidth}{XX}
\toprule
\textbf{Mobile-LiDAR idea} & \textbf{Classical surveying equivalent}\\ \midrule
Time-of-flight range & EDM distance\\
Polar-to-Cartesian conversion & Total-station angle + distance $\to$ coordinates\\
Pose (position + orientation) & Instrument station + backsight orientation\\
IMU dead reckoning & An open traverse --- error accumulates\\
Loop closure & Closing a traverse\\
Factor-graph optimization & Least-squares network adjustment\\
Control points vs checkpoints & Constrained points vs independent check shots\\
Error budget (root-sum-square) & Traverse / level-network design\\
\bottomrule
\end{tabularx}
\end{center}

\section{What you will be able to do at the end}
\begin{itemize}
\item Compute a range from a time of flight and explain why we divide by two.
\item Turn a range and two angles into $X,Y,Z$ --- and then into project coordinates using lever arm, boresight, attitude and trajectory.
\item Predict how fast an IMU of a given grade drifts, and how far a timing or boresight error moves a point.
\item Explain, with numbers, what an extended Kalman filter, ICP, and a factor graph actually do.
\item Read the LIO-SAM paper's equations and know what every symbol means.
\item Build an error budget, compute RMSE from checkpoints, and report accuracy in the language of current standards.
\end{itemize}

\begin{tryit}
\textbf{1.1} (Simple) If you clap and hear the echo after 1 second, and sound travels about 343 meters each second, how far away is the wall? (Hint: the sound went there \emph{and} back.)\\
\textbf{1.2} (Engineer) List the five ingredients (besides the laser) needed to turn a moving scanner's ranges into a survey-grade map.
\end{tryit}

\sources{Shan, T.\ \emph{et al.} (2020), LIO-SAM, IEEE/RSJ IROS, arXiv:2007.00258; \emph{Mobile LiDAR Basics for Surveying and Mapping} seminar (2026); \emph{LiDAR Mobile Mapping with SLAM} course outline (2026).}
